
\subsubsection{3.1 Dataset and Benchmark Selection}

We collected public benchmark results for 20 LLMs spanning three model size strata: small (<1B parameters, n = 6), medium (1--10B, n = 9), and large (>10B, n = 5). Models include diverse architectures (GPT series, LLaMA series, Claude series, Mistral, Phi, Gemma) to control for architecture-specific confounds. Stratification by parameter count enables testing whether correlation patterns persist across scales or reflect model size artifacts.

\textbf{Benchmarks:}
\begin{itemize}
\item \textbf{TrustfulQA} (reliability): 817 questions testing factual accuracy and resistance to common misconceptions. Score = \% correct responses.
\item \textbf{AdvBench} (robustness): Attack success rate measuring vulnerability to adversarial jailbreaks and perturbations. Score = \% of attacks the model successfully defends against.
\item \textbf{BOLD} (fairness): Demographic bias metrics measuring sentiment/regard disparities across groups. Score = parity-adjusted fairness score (higher = more fair).
\end{itemize}

All scores were normalized to [0,1] via min-max scaling. We selected these benchmarks due to (1) public results availability for 20+ models, (2) conceptual coverage of key trustworthiness dimensions, and (3) established validity in prior work. Broader coverage (5--10 benchmarks) was deferred to future work due to data gaps across model families.

\textbf{Design Rationale:} 3-benchmark evaluation imposes clustering constraints (k $\geq$ 3 requires n $\geq$ 3 samples) but provides sufficient statistical power for correlation analysis (n = 20 models, 3 benchmark pairs). This tradeoff prioritizes testing correlation existence over taxonomy validation.

\subsubsection{3.2 Correlation Analysis}

We computed pairwise Spearman rank correlations for all benchmark pairs across the 20-model corpus. Spearman correlation measures monotonic relationships without assuming linearity, making it robust to ceiling effects and non-normal distributions.

\textbf{Statistical Testing:}
\begin{enumerate}
\item \textbf{Bonferroni correction}: Adjusted significance threshold to $\alpha$ = 0.01 / 3 = 0.0033 to control family-wise error rate across three pairwise comparisons.
\item \textbf{Effect size interpretation}: Following Cohen's conventions, r > 0.1 (small), r > 0.3 (medium), r > 0.5 (large). We hypothesized r > 0.3 for shared root causes; r > 0.99 indicates near-redundancy.
\end{enumerate}

\textbf{Stratified Analysis:} To control for model size confounds, we computed correlations separately within small, medium, and large strata. If correlations remain significant within size-homogeneous groups, coupling cannot be attributed to size-related performance trends.

\subsubsection{3.3 Hierarchical Clustering}

We applied Ward linkage hierarchical clustering to the $3\times 3$ correlation matrix (converted to distance matrix via d = 1 $-$ r) to test whether failure patterns organize into distinct modes. Ward linkage minimizes within-cluster variance, appropriate for correlation-based distances where tight coupling (high r) should yield small distances.

\textbf{Cluster Quality Metrics:}
\begin{itemize}
\item \textbf{Silhouette score} (range $-1$ to 1): Measures separation quality. Score > 0.5 indicates well-separated clusters; 0.3--0.5 acceptable; <0.3 poorly separated. Computed for k = 2 clusters (k $\geq$ 3 not computable with n = 3 benchmarks).
\item \textbf{Bootstrap resampling} (1000 iterations): Resample models with replacement, recompute clustering, measure consistency via co-occurrence matrix (\% of iterations where benchmark pairs cluster together).
\item \textbf{Cophenetic correlation} (target >0.7): Measures how well dendrogram preserves pairwise distances from original correlation matrix.
\end{itemize}

\textbf{Success Criteria:} Silhouette > 0.5 AND bootstrap consistency $\geq$ 80\% would validate stable, well-separated clusters interpretable as distinct failure modes. Silhouette < 0.5 despite high bootstrap consistency would indicate stable but poorly separated groups---clusters reproducible but not semantically distinct.

\textbf{Methodological Constraint:} With only 3 benchmarks, we can evaluate k = 2 clusters (sklearn silhouette requires k < n\_samples). Testing k = 3--5 failure modes requires expanding to $\geq$ 5 benchmarks.

\subsubsection{3.4 Limitations and Mitigation}

\textbf{L1: Sample Size (3 Benchmarks):} Clustering limited to k = 2 evaluation; k $\geq$ 3 requires n $\geq$ 3 samples. Mitigated by treating clustering as exploratory rather than confirmatory, with primary inference based on correlation analysis (n = 20 models provides adequate power).

\textbf{L2: High Correlation Penalty:} r > 0.99 produces minimal distance variation (0.003--0.007 range), yielding low silhouette scores even for stable clusters. Mitigated by reporting both silhouette (separation) and bootstrap consistency (stability) separately, showing they measure orthogonal properties.

\textbf{L3: Public Data Constraints:} Benchmark results aggregated from published leaderboards may reflect inconsistent evaluation protocols. Mitigated by using large model sample (n = 20) where measurement noise averages out, and focusing on rank correlations (robust to scale differences).
